\section*{Appendix E: Path-Ordered Holonomy and Loop Transport}
\addcontentsline{toc}{section}{Appendix E: Path-Ordered Holonomy and Loop Transport}
\label{app:holonomy}

\begin{definition}[Phase-Space Holonomy Integral]
Let $\gamma : [0, 1] \to T^*\M$ be a closed trajectory loop in canonical position-momentum coordinates $(\mathbf{q}, \mathbf{p})$ bounding a two-dimensional surface $S$. By Stokes' theorem, the accumulated holonomy displacement is given by the path-ordered line integral:
\begin{equation}
  \mathrm{Hol}(\gamma) = \oint_\gamma \mathbf{p} \cdot d\mathbf{q} = \iint_S \Omega
\end{equation}
where $\Omega$ denotes the symplectic curvature two-form \cite{earn2001exact}. Under un-dithered quantization noise or out-of-order operator execution, $\mathrm{Hol}(\gamma) \neq 0$ at mechanical rest, manifesting empirically as recurrent token-looping hallucinations \cite{anonymous2026geometry}.
\end{definition}

\begin{proposition}[Unitary Neutral Trace Invariance]
Let $\tau = [t_1, t_2, \dots, t_k]$ be a discrete sequence of balanced-ternary state transitions over $\mathcal{T} = \{T^-, T^0, T^+\}$. If every transition in $\tau$ evaluates to the neutral identity element ($t_i = T^0$), sequential multiplicative folding preserves the unitary anchor identically:
\begin{equation}
  \mathrm{foldl}(\otimes, T^+, \tau) = T^+.
\end{equation}
\end{proposition}

\begin{proof}
Established via structural list induction in linear time $\mathcal{O}(k)$ over MathComp sequence predicates (\texttt{allP}) in \texttt{Holonomy.v}, eliminating exponential case-splitting overhead \cite{barrett2024smt}.
\end{proof}
